%%%%%%%%%%%%%%%%%%%%%%%%%%%%%%%%%%%%%%%%%
% Professional Formal Letter
% LaTeX Template
% Version 2.0 (12/2/17)
%
% This template originates from:
% http://www.LaTeXTemplates.com
%
% Authors:
% Brian Moses
% Vel (vel@LaTeXTemplates.com)
%
% License:
% CC BY-NC-SA 3.0 (http://creativecommons.org/licenses/by-nc-sa/3.0/)
%
%%%%%%%%%%%%%%%%%%%%%%%%%%%%%%%%%%%%%%%%%


%----------------------------------------------------------------------------------------
%	PACKAGES AND OTHER DOCUMENT CONFIGURATIONS
%----------------------------------------------------------------------------------------

\documentclass[11pt, a4paper]{letter} % Set the font size (10pt, 11pt and 12pt) and paper size (letterpaper, a4paper, etc)

%\graphicspath{{./fig/}} %fig path NOT WORKING
\usepackage{hyperref}


%%%%%%%%%%%%%%%%%%%%%%%%%%%%%%%%%%%%%%%%%
% Professional Formal Letter
% Structure Specification File
% Version 1.0 (12/2/17)
%
% This file originates from:
% http://www.LaTeXTemplates.com
%
% Authors:
% Brian Moses
% Vel (vel@LaTeXTemplates.com)
%
% License:
% CC BY-NC-SA 3.0 (http://creativecommons.org/licenses/by-nc-sa/3.0/)
%
%%%%%%%%%%%%%%%%%%%%%%%%%%%%%%%%%%%%%%%%%

%----------------------------------------------------------------------------------------
%	PACKAGES AND OTHER DOCUMENT CONFIGURATIONS
%----------------------------------------------------------------------------------------

\usepackage{graphicx} % Required for including pictures

\usepackage[T1]{fontenc} % Output font encoding for international characters
\usepackage[utf8]{inputenc} % Required for inputting international characters
\usepackage{gfsdidot} % Use the GFS Didot font: http://www.tug.dk/FontCatalogue/gfsdidot/
\usepackage{microtype} % Improves typography

\pagestyle{empty} % Suppress headers and footers

\setlength\parindent{1cm} % Paragraph indentation

% Create a new command for the horizontal rule in the document which allows thickness specification
\makeatletter
\newcommand{\vhrulefill}[1]{\leavevmode\leaders\hrule\@height#1\hfill \kern\z@}
\makeatother

%----------------------------------------------------------------------------------------
%	DOCUMENT MARGINS
%----------------------------------------------------------------------------------------

\usepackage{geometry} % Required for adjusting page dimensions

\geometry{
	top=1cm, % Top margin
	bottom=1.5cm, % Bottom margin
	left=3cm, % Left margin
	right=3cm, % Right margin
	%showframe, % Uncomment to show how the type block is set on the page
}

%----------------------------------------------------------------------------------------
%	DEFINE CUSTOM COMMANDS
%----------------------------------------------------------------------------------------

\newcommand{\logo}[1]{\renewcommand{\logo}{#1}}

\newcommand{\Who}[1]{\renewcommand{\Who}{#1}}
\newcommand{\Title}[1]{\renewcommand{\Title}{#1}}

\newcommand{\headerlineone}[1]{\renewcommand{\headerlineone}{#1}}
\newcommand{\headerlinetwo}[1]{\renewcommand{\headerlinetwo}{#1}}

\newcommand{\authordetails}[1]{\renewcommand{\authordetails}{#1}}

%----------------------------------------------------------------------------------------
%	AUTHOR DETAILS STRUCTURE
%----------------------------------------------------------------------------------------

\newcommand{\authordetailsblock}{
	\hspace{\fill} % Move the author details to the far right
	\parbox[t]{0.48\textwidth}{ % Box holding the author details; width value specifies where it starts and ends, increase to move details left
		\footnotesize % Use a smaller font size for the details
		\Who\\ % Author name
		\textit{\authordetails} % The author details text, all italicised
	}
}

%----------------------------------------------------------------------------------------
%	HEADER STRUCTURE
%----------------------------------------------------------------------------------------

\address{
	\includegraphics[width=1in]{\logo} % Include the logo of author institution
	\hspace{0.82\textwidth} % Position of the institution logo, increase to move left, decrease to move right
	\vskip -0.1\textheight~\\ % Position of the large header text in relation to the institution logo, increase to move down, decrease to move up
	\Large\hspace{0.2\textwidth}\headerlineone\hfill ~\\[0.006\textheight] % First line of institution name, adjust hspace if your logo is wide
	\hspace{0.2\textwidth}\headerlinetwo\hfill \normalsize % Second line of institution name, adjust hspace if your logo is wide
	\makebox[0ex][r]{\textbf{\Who\Title}}\hspace{0.01\textwidth} % Print author name and title with a little whitespace to the right
	~\\[-0.01\textheight] % Reduce the whitespace above the horizontal rule
	\hspace{0.2\textwidth}\vhrulefill{1pt} \\ % Horizontal rule, adjust hspace if your logo is wide and \vhrulefill for the thickness of the rule
	\authordetailsblock % Include the letter author's details on the right side of the page under the horizontal rule
	\hspace{-0.25\textwidth} % Horizontal position of the author details block, increase to move left, decrease to move right
	\vspace{-0.1\textheight} % Move the date and letter content up for a more compact look
}

%----------------------------------------------------------------------------------------
%	COMPOSE THE ENTIRE HEADER
%----------------------------------------------------------------------------------------

\renewcommand{\opening}[1]{
	{\centering\fromaddress\vspace{0.05\textheight} \\ % Print the header and from address here, add whitespace to move date down
	\hspace*{\longindentation}\today\hspace*{\fill}\par} % Print today's date, remove \today to not display it
	{\raggedright \toname \\ \toaddress \par} % Print the to name and address
	\vspace{1cm} % White space after the to address
	\noindent #1 % Print the opening line
}

%----------------------------------------------------------------------------------------
%	SIGNATURE STRUCTURE
%----------------------------------------------------------------------------------------

\signature{\Who\Title} % The signature is a combination of the author's name and title

\renewcommand{\closing}[1]{
	\vspace{2.5mm} % Some whitespace after the letter content and before the signature
	\noindent % Stop paragraph indentation
	\hspace*{\longindentation} % Move the signature right to the value of \longindentation
	\parbox{\indentedwidth}{
		\raggedright
		#1 % Print the signature text
		\vskip 1.65cm % Whitespace between the closing text and author's name for a physical signature
		\fromsig % Prints the value of \signature{}, i.e. author name and title
	}
}
 % Include the file that specifies the document structure

\longindentation=0pt % Un-commenting this line will push the closing "Sincerely," and date to the left of the page

%----------------------------------------------------------------------------------------
%	YOUR INFORMATION
%----------------------------------------------------------------------------------------

%\Who{} % Your name
%\Title{} % Your title, leave blank for no title

\authordetails{%
%	Département de mathematiques\\ et de statistique\\ % Your department/institution
	Samuel Perreault\\
	623 rue Du Moulin\\ % Your address
	Mont-Tremblant, Qc, J8E 2V1\\ % Your city, zip code, country, etc
	Email:  communications@dotlayer.org\\ % Your email address
	Phone: (819) 425-0303\\ % Your phone number
	URL: \href{https://www.dotlayer.org/}{dotlayer.org} % Your URL
}

%----------------------------------------------------------------------------------------
%	HEADER CONTENTS
%----------------------------------------------------------------------------------------

\logo{logo.png} % Logo filename, your logo should have square dimensions (i.e. roughly the same width and height), if it does not, you will need to adjust spacing within the HEADER STRUCTURE block in structure.tex (read the comments carefully!)

\headerlineone{} % Top header line, leave blank if you only want the bottom line

\headerlinetwo{.Layer} % Bottom header line

%----------------------------------------------------------------------------------------

\begin{document}

%----------------------------------------------------------------------------------------
%	TO ADDRESS
%----------------------------------------------------------------------------------------

\begin{letter}{
	Prof. Mary Thompson\\
	Chair of the Fundraising Committee\\
	Statistical Society of Canada
	
	\bigskip
	\textbf{Subject: education and outreach proposal}%
}

%----------------------------------------------------------------------------------------
%	LETTER CONTENT
%----------------------------------------------------------------------------------------


\opening{Dear Prof. Thompson,}

%\emph{CALL: This is a new call for proposals for advancing the goals of the SSC to promote statistics and probability to students and educators, or to help students and young researchers develop careers in the statistical sciences.}
%
%\emph{The SSC Board will award between 500 and 5,000 for the best proposals received by January 31, 2019. Preference will be given to proposals that have a broad reach, e.g. are regional or national rather than local, and that are able to obtain additional support. The number of awards will depend on the funds available. If a proposal involves electronic services, proposers are encouraged to contact the Electronic Services Manager of the SSC for advice.}

I am writing to you on behalf of \textit{.Layer} (pronounced ``dot layer''), a community created for promoting collaboration and knowledge sharing in the field of data science. As an organization in its creation process, we were very pleased to learn about the recent call for education and outreach proposals of the Statistical Society of Canada. I will try, in this letter, to clearly describe who we are, what our mission is, and how we intend to put our words into action. We hope that you will find our proposal convincing enough to grant us your support.

%\pagebreak
%
%\emph{QUESTION: Your idea and a rationale, explaining how it fits with the mission of the SSC:}
%   
%\emph{INFO: The Statement of Purpose of the SSC is to advance the knowledge and education of Canadians concerning statistical sciences and related fields.}

%As a statistics enthusiast, I rarely doubt my decision of pursuing graduate studies in statistics. However, it is hard to deny, as you may have experienced yourself, that such enterprise sometimes comes with its lot of solitude. To my great surprise,

\bigskip
\noindent \textbf{Who we are}

Last Summer, some of us organized a week-end of presentations and activities related to data science. Friends of friends, fifteen in total, then met for an epic brainstorm. Out of it emerged \emph{.Layer}, a hub consisting of more or less 25 dedicated students and/or practitioners in the fields of statistics, actuarial science and computer science, (mostly) residing in either Montreal or Quebec City. We create, through various initiatives, both virtual and physical spaces where people can meet, learn and share. Members responsible for the present proposal are (* = SSC member):
\begin{enumerate}
	\item[] \textbf{Jean-Thomas Baillargeon, MSc FSA, CERA}\\
	\quad PhD student at Laval University (computer science)\\
	\quad Lecturer at Laval University\\
	\quad Former actuarial analyst at Industrial Alliance (6 years)\\
	\quad Former data product manager at Xpertsea Solution (2 years)
	\item[] \textbf{David Beauchemin, BSc}\\
	\quad Masters student at Laval University (computer science)
	\item[] \textbf{Christopher Blier-Wong, MSc}\\
	\quad PhD student at Laval University (actuarial science)\\
	\quad Masters student at Laval University (computer science)
	%\\\quad Member of university reseach labs: ACT\&RISK, QUANTACT and BDRC
	\item[] \textbf{Antoine Buteau, BCom}\\
	\quad TMT Management Consultant at The PNR\\
	\quad Former solution specialist at Coveo (2.5 years)\\
	\quad Former board advisor at Caisse Desjardins -- Chaudière et Etchemin (2 years)
	\item[] \textbf{Stéphane Caron*, BSc}\\
	\quad Actuarial analyst at Intact Insurance\\
	\quad Masters student at Laval University (mathematics and statistics)
	\item[] \textbf{Marie-Pier Côté*, PhD, FSA}\\
	\quad Assistant professor, École d'actuariat, Laval University\\
	\quad Teaching chair on big data analysis for actuarial science --- Intact\\
	\quad Member of the \emph{New Investigators Committee} of the \emph{SSC}
	\item[] \textbf{Nicolas Garneau, BSc}\\
	\quad PhD Student at Laval University (Computer Science)\\
	Lecturer at Laval University, Computer Science, Deep Learning\\
	Deep Learning Trainer, Big Data Research Center
	\item[] \textbf{Samuel Perreault*, MSc}\\
	\quad PhD student at Laval University (mathematics and statistics)\\
	\quad Research \& innovation analyst at The Cooperators Insurance
\end{enumerate}

\noindent Three of us, Marie-Pier Côté, Christopher Blier-Wong and Samuel Perreault presented at the 2018 SSC meeting.

\noindent So far, we have been working mostly on the two following initiatives.
\begin{enumerate}
	\item[(a)] \emph{Meetup Machine Learning Québec}: This project was founded by Nicolas Garneau. As the name suggests, it is dedicated to the organization of meetups and conferences related to machine learning in Quebec City. Before \emph{.Layer} even existed, \emph{Meetup} had already organized three events in 2017-2018 (at Laval University/XpertSea's offices/Coveo's offices), attracting more than 100 people in total. Since then, Nicolas joined forces with many others to organize events of even greater scope. Concretely, we organized both a conference (Laval University) and a hackathon (Bloc Solutions' offices) on statistics and machine learning applied to casualty insurance; as well as a workshop on natural language processing (Laval University). These events attracted an approximate total of 200 people. They were paid for (rooms, food, beverages, web services for GPUs, etc.) by our sponsors: Intact Insurance, Co-operators Insurance, Bloc Solutions and, from Laval University, the Big Data Research Center (CRDM), the Chaire d'actuariat and the service de placements (SPLA).
	\item[(b)] \emph{.Layer blog}: During the founding week-end of \emph{.Layer}, we agreed that a blog would be great as a first project to collaborate on, as well as a convenient way of reporting our activities and sharing our ideas. Though great devotion: the blog is already up and running at \href{https://www.dotlayer.org/}{dotlayer.org}. In order to reach as many people as possible, we have recently made the website bilingual (en/fr).
\end{enumerate}

The next big step for us is the registration of a nonprofit organization for administering the community's activities. Here is a sketch of the structure we intend to put in place.

\begin{enumerate}
	\item[] \textbf{Board of directors}: The board consists of five people: an executive president, a secretary, a treasurer, a vice-president assigned to communications and a vice-president assigned to external affairs.
	\item[] \textbf{Members}: In addition to contributing and participating to the organization's projects and events, members have the right to vote during the general assemblies. One must first be accepted by the current members in order to gain its member status.
	\item[] \textbf{Sympathizers}: Any other individual involved in the community's projects and events. They can share their opinion, but have no right to vote during assemblies.
\end{enumerate}

This simple structure allows the organization to be as inclusive as possible while assuring that funds and efforts are spent on activities that are aligned with its mission. Any sympathizer minimally involved in the community, and demonstrating that he or she shares its inclusive values, should expect to be accepted as a member if requested.

\pagebreak

\noindent \textbf{Three distinct proposals}

Since \emph{.Layer} is not a project in itself, but a hub for projects, we have decided to split our proposal in three distinct parts. The first concerns \emph{.Layer} in general: we seek help for covering the costs related the organization of a healthy and inclusive community. The two others are related to web-based initiatives and physical events respectively. Details follow.

\begin{enumerate}
	\item \textbf{Structuring the community}
	
	\quad We are lucky to have among us a professional graphic designer, Jean-Christophe Yelle, who created our logo at no cost. However, we believe that professional services not directly related to data science (like logo design, t-shirt design, social media setup, and so on) should be fairly remunerated. In the budget provided with the letter, we also detail recent and future expenses related to the registration of \emph{.Layer} as a nonprofit organization.
	
	\bigskip
	\emph{Member in charge:} Stéphane Caron\\
	
	\item \textbf{Web-based material}
	
	\quad We count on the \emph{.Layer blog} to get a reach far beyond Quebec City and Montreal. We are committed to providing web material of great quality, and we think that it begins with a great use of the french and english languages. However, producing relevant material is a time-consuming enterprise and the translation of the articles we judge important can quickly become a burden. To make sure everyone spends his or her time on data science rather than translation, we would like get revision services from a professional translator. People willing to have their article in both languages would then only have to provide a draft of a translation. An overall revision of the website would be warmly welcomed too.
	
	\quad Web-based material is not restricted to the blog. Some of us are currently working on a podcast named \emph{OpenLayer}. In video clips of about 45-60 minutes, David Beauchemin would meet different actors of the data science community to discuss the journey that took them where they are today and what their day-to-day looks like. The objective is to shed light on the job market, a hidden layer for aspiring data scientists. David could for example meet a Ph.D. student in statistics or artificial intelligence or a data analyst at some insurance company. The concept is influenced by podcasts \emph{à la} Joe Rogan and ``Mike Ward sous écoute''. In each episode, we will discover someone and his/her environment through a casual conversation.
	
	\bigskip
	\emph{Members in charge:} Antoine Buteau and David Beauchemin\\
	
	\item \textbf{Educational event on statistical learning}.
	
	\quad More \emph{Meetup Machine Learning Québec} events are coming. An introduction to the foundations of machine learning, followed by a workshop on computer vision are offered on February 7 and 9 respectively. We are always seeking funds for organizing more events like these. In particular, we would like to organize one focusing on so-called statistical learning (we are big fans of the book \emph{The Elements of Statistical Learning}). Your support, with that of other sponsors from the statistical community, would help us pay for rooms, food and beverages as well as keeping the price as low as possible for attendees. This would be a great opportunity for students to learn about the frontier between statistics and machine learning, in addition to being great opportunity for meeting professionals of various domains in which statistics is put to use.

			
	\bigskip
	\emph{Members in charge:} Nicolas Garneau and Christopher Blier-Wong\\
	
\end{enumerate}


\bigskip
\noindent \textbf{Administration of the funds}

Administering the funds and reporting how they are used is the \emph{raison d'\^{e}tre} of the nonprofit entity tied to .Layer. Until its official incorporation, the members responsible for the proposal will take the responsability of providing such report to the \emph{Statistical Society of Canada}. We provided a budget and a timeline for each of the three proposals detailed in this letter.

\bigskip

We strongly believe that the goals of \emph{.Layer} are aligned with the mission of the \emph{SSC}. Through our workshops of varying levels and our educational blog posts, we help to broadcast the importance of sound statistical methods for data analysis. The blog acts as a platform for sharing ideas and reinforcing the sense of community between all individuals interested in data science and related fields. The meetup events and the workshops are also great opportunities to foster exchanges between theoreticians and practitioners. On behalf of everyone involved in \emph{.Layer}, I would like to thank you for considering our proposal. We wish to have convinced you that our projects are worthy of your support.

\closing{Sincerely,}

\noindent Samuel Perreault

%----------------------------------------------------------------------------------------
%	OPTIONAL FOOTNOTE
%----------------------------------------------------------------------------------------

% Uncomment the 4 lines below to print a footnote with custom text
%\def\thefootnote{}
%\def\footnoterule{\hrule}
%\footnotetext{\hspace*{\fill}{\footnotesize\em Footnote text}}
%\def\thefootnote{\arabic{footnote}}

%----------------------------------------------------------------------------------------

\end{letter}

\end{document}
